% ═══ 由 环境/pdf/latex/md到LaTeX.py 生成 ═══
% 中文字体：思源宋体（子集，SIL OFL 1.1）；拉丁/数学：Latin Modern（随 TeX Live 分发）
% 编译：xelatex（本仓库自带 wasm 版 XeTeX 引擎，见 环境/pdf/latex/编译LaTeX.mjs）
\usepackage[T1]{fontenc}      % 否则 ASCII 的 < > | 在 OT1 下排成 ¡ ¿ —
\usepackage{lmodern}          % 必需：数学的物理字模（否则 xdvipdfmx 拒绝出 PDF）
\usepackage{amsmath,amssymb,mathtools}
\usepackage{booktabs,longtable,array,multirow}
\usepackage{graphicx}
\usepackage{xcolor}
\usepackage{listings}
% 不用 hyperref：新版 hyperref 强制依赖 bookmark.sty，而 bundle 里没有。
% 链接命令用降级定义，保证正文里的 \href / \url 不会报错。
\providecommand{\href}[2]{#2}
\providecommand{\url}[1]{\texttt{#1}}
\providecommand{\hypertarget}[2]{#2}
\providecommand{\hyperlink}[2]{#2}
\providecommand{\urlstyle}[1]{}
\providecommand{\texorpdfstring}[2]{#1}   % hyperref 的命令，剥离后仍被模板/标题用到
\providecommand{\phantomsection}{}
\providecommand{\pdfstringdef}[2]{}
\providecommand{\hypersetup}[1]{}

% ── 三套字体：中文 / 符号 / emoji（按字符 cmap 自动选择）──
\font\zhfont="[NotoSerifSC-Regular.ttf]:script=hani" at 10.5pt
\font\symfont="[DejaVuSans.ttf]" at 10.5pt
\font\emofont="[NotoEmoji-Regular.ttf]" at 10.5pt
\font\zhmonofont="[NotoSansSC-Regular.ttf]" at 9pt
% 用法：\zh{中文} / \sym{→} / \emo{🦙}
% 注意：必须**带参数**。写成 \def\zh{{\zhfont}} 是错的——字体赋值只在小群里有效，
% 那个小群随 \zh 立刻结束，于是中文仍然用拉丁字体排，报 Missing character。
\def\zh#1{{\zhfont #1}}
\def\sym#1{{\symfont #1}}
\def\emo#1{{\emofont #1}}
% 含汉字的行内代码：listings 的 \lstinline 遇到汉字必炸（`\lst@arg ->判`
% 未定义控制字，catcode 设 12 也无效），所以直接自己排。
\def\codezh#1{{\zhmonofont #1}}

\def\zhglue{\hskip0pt plus .06em minus .01em}
\linespread{1.12}
\setlength{\parindent}{0pt}
\setlength{\parskip}{0.45em}

% ── 代码块：等宽 + 中文字体（listings 默认字体不含汉字）──
\lstset{
  basicstyle=\zhmonofont,
  breaklines=true,
  breakatwhitespace=false,
  columns=fullflexible,
  keepspaces=true,
  showstringspaces=false,
  frame=single,
  framesep=4pt,
  rulecolor=\color{gray!45},
  backgroundcolor=\color{gray!7},
  xleftmargin=6pt,
  extendedchars=true,
  tabsize=2,
  % 代码块里的 emoji：NotoSansSC 没有这些字形（listings 不做字符级回退）。
  % 用 escapeinside 开口子，转换器把 emoji 逐字包成 (*@{\emo{⭐}}@*)。
  % 试过 literate={⭐}{{\emofont ⭐}}1，会把 ASCII 字母也吞掉并报
  % `Improper alphabetic constant`，不可用。
  escapeinside={(*@}{@*)},
}
\renewcommand{\lstlistingname}{代码}

% ── 表格线宽（booktabs 风格）──
\renewcommand{\arraystretch}{1.25}

% ── 标题样式 ──
\usepackage{titlesec}
\titleformat{\section}{\Large\bfseries}{\thesection}{0.6em}{}
\titlespacing*{\section}{0pt}{1.1em}{0.5em}
